\documentclass[10pt,a4paper]{amsart}
\usepackage{amsmath,amssymb,amsthm,graphicx}
\usepackage{verbatim}
\newenvironment{reference}{\comment}{\endcomment}
%\newenvironment{reference}{}{}
\newenvironment{slogan}{\comment}{\endcomment}
\newenvironment{history}{\comment}{\endcomment}
% For commutative diagrams we use Xy-pic
\usepackage[all]{xy}

% We use 2cell for 2-commutative diagrams.
\xyoption{2cell}
% Theorem environments.
%
\theoremstyle{plain}
\newtheorem{theorem}[subsection]{Theorem}
\newtheorem{proposition}[subsection]{Proposition}
\newtheorem{lemma}[subsection]{Lemma}

\theoremstyle{definition}
\newtheorem{definition}[subsection]{Definition}
\newtheorem{example}[subsection]{Example}
\newtheorem{exercise}[subsection]{Exercise}
\newtheorem{situation}[subsection]{Situation}

\theoremstyle{remark}
\newtheorem{remark}[subsection]{Remark}
\newtheorem{remarks}[subsection]{Remarks}

\numberwithin{equation}{subsection}

% Macros
%
\def\lim{\mathop{\mathrm{lim}}\nolimits}
\def\colim{\mathop{\mathrm{colim}}\nolimits}
\def\Spec{\mathop{\mathrm{Spec}}}
\def\Hom{\mathop{\mathrm{Hom}}\nolimits}
\def\Ext{\mathop{\mathrm{Ext}}\nolimits}
\def\SheafHom{\mathop{\mathcal{H}\!\mathit{om}}\nolimits}
\def\SheafExt{\mathop{\mathcal{E}\!\mathit{xt}}\nolimits}
\def\Sch{\mathit{Sch}}
\def\Mor{\mathop{\mathrm{Mor}}\nolimits}
\def\Ob{\mathop{\mathrm{Ob}}\nolimits}
\def\Sh{\mathop{\mathit{Sh}}\nolimits}
\def\NL{\mathop{N\!L}\nolimits}
\def\CH{\mathop{\mathrm{CH}}\nolimits}
\def\proetale{{pro\text{-}\acute{e}tale}}
\def\etale{{\acute{e}tale}}
\def\QCoh{\mathit{QCoh}}
\def\Ker{\mathop{\mathrm{Ker}}}
\def\Im{\mathop{\mathrm{Im}}}
\def\Coker{\mathop{\mathrm{Coker}}}
\def\Coim{\mathop{\mathrm{Coim}}}

% Boxtimes
%
\DeclareMathSymbol{\boxtimes}{\mathbin}{AMSa}{"02}

%
% Macros for moduli stacks/spaces
%
\def\QCohstack{\mathcal{QC}\!\mathit{oh}}
\def\Cohstack{\mathcal{C}\!\mathit{oh}}
\def\Spacesstack{\mathcal{S}\!\mathit{paces}}
\def\Quotfunctor{\mathrm{Quot}}
\def\Hilbfunctor{\mathrm{Hilb}}
\def\Curvesstack{\mathcal{C}\!\mathit{urves}}
\def\Polarizedstack{\mathcal{P}\!\mathit{olarized}}
\def\Complexesstack{\mathcal{C}\!\mathit{omplexes}}
% \Pic is the operator that assigns to X its picard group, usage \Pic(X)
% \Picardstack_{X/B} denotes the Picard stack of X over B
% \Picardfunctor_{X/B} denotes the Picard functor of X over B
\def\Pic{\mathop{\mathrm{Pic}}\nolimits}
\def\Picardstack{\mathcal{P}\!\mathit{ic}}
\def\Picardfunctor{\mathrm{Pic}}
\def\Deformationcategory{\mathcal{D}\!\mathit{ef}}
\title{Uniform vanishing of a Tor map in the degree of an upper bound for the number of generators of the ideal}
\author{The Stacks Project Authors}
\date{stacks-project:local-cohomology.tex:local-snapshot-20261012; GNU FDL 1.2 or later}
\begin{document}
\maketitle
\noindent\textit{Editorial note.} All delivered mathematical text is the author's own, copied from the pinned Stacks Project chapter source sources/stacks-local-cohomology/pinned-original-source.tex; no statement or proof step was rewritten, reordered or removed.\par\medskip
\noindent\textit{Editorial note.} Preamble provenance: the Stacks Project keeps one single shared preamble for every chapter; this card takes it directly from the staged sources/stacks-local-cohomology/preamble.tex as three verbatim spans (lines 8-12, 14-18 and 154-217), so no derived file was written and nothing was added or redefined. The wrapper class and the base packages are supplied by the assembler.\par\medskip
\noindent\textit{Editorial note.} Bibliography provenance: neither the counted statement nor its proof contains a citation (this lane rescans the whole span before writing the card), so the card carries no bibliography and no entry was reconstructed.\par\medskip
\noindent\textit{Editorial note.} Selection provenance: this card was opened under the file-based control contract (CONTROL.md), with the mechanical pre-card gate run on its statement and proof interval before the card was written.\par\medskip
\noindent\textit{Allowed standard prerequisites (editorial).}\begin{itemize}
\item The annihilation of the Tor modules above the number of generators of the ideal by a fixed power of the ideal, uniformly in the module and the exponent. Author lemma lemma-annihilates-tors, whose statement is reproduced verbatim as a context span of this card (source lines 4668-4675). Supplies the single constant c\_1 used in the definition of the invariant of the module and in the diagram of the induction step.
\item The vanishing of the Tor maps for a module annihilated by a fixed power of the ideal, again uniformly in the module. Author lemma lemma-maps-zero-fixed-torsion, whose statement is reproduced verbatim as a context span of this card (source lines 4510-4519). Supplies the constants that give the base case of the induction.
\end{itemize}\medskip
\begin{lemma}
\label{lemma-tor-maps-vanish}
Let $I$ be an ideal of a Noetherian ring $A$. Let $t \geq 0$
be an upper bound on the number of generators of $I$.
There exist $N, c \geq 0$ such that the maps
$$
\text{Tor}_{t + 1}^A(M, A/I^n) \to \text{Tor}_{t + 1}^A(M, A/I^{n - c})
$$
are zero for any $A$-module $M$ and all $n \geq N$.
\end{lemma}

\begin{lemma}
\label{lemma-annihilates-tors}
In the situation above, let $t$ be an upper bound on the number of
generators for $I$. There exists an integer $c = c(A, I) \geq 0$
such that for any $A$-module $M$ the tor modules
$\text{Tor}_i^A(M, A/I^q)$ are annihilated by $I^c$ for $i > t$
and all $q \geq 0$.
\end{lemma}

\begin{lemma}
\label{lemma-maps-zero-fixed-torsion}
Let $I$ be an ideal of a Noetherian ring $A$. For every $m \geq 0$
and $i > 0$ there exist a $c = c(A, I, m, i) \geq 0$ such that
for every $A$-module $M$ annihilated by $I^m$ the map
$$
\text{Tor}^A_i(M, A/I^n) \to \text{Tor}^A_i(M, A/I^{n - c})
$$
is zero for all $n \geq c$.
\end{lemma}

\begin{proof}
Let $c_1$ be the constant found in Lemma \ref{lemma-annihilates-tors}.
Please keep in mind that this constant $c_1$ works for $\text{Tor}_i$ for all
$i > t$ simultaneously.

\medskip\noindent
Say $I = (a_1, \ldots, a_t)$. For an $A$-module $M$ we set
$$
\ell(M) =
\#\{i \mid 1 \leq i \leq t,\ a_i^{c_1}\text{ is zero on }M\}
$$
This is an element of $\{0, 1, \ldots, t\}$.
We will prove by descending induction on $0 \leq s \leq t$
the following statement $H_s$: there exist $N, c \geq 0$ such that
for every module $M$ with $\ell(M) \geq s$ the maps
$$
\text{Tor}_{t + 1 + i}^A(M, A/I^n) \to \text{Tor}_{t + 1 + i}^A(M, A/I^{n - c})
$$
are zero for $i = 0, \ldots, s$ for all $n \geq N$.

\medskip\noindent
Base case: $s = t$. If $\ell(M) = t$, then $M$ is annihilated by
$(a_1^{c_1}, \ldots, a_t^{c_1}\}$ and hence by
$I^{t(c_1 - 1) + 1}$. We conclude from
Lemma \ref{lemma-maps-zero-fixed-torsion}
that $H_t$ holds by taking $c = N$ to be the maximum of the integers
$c(A, I, t(c_1 - 1) + 1, t + 1), \ldots, c(A, I, t(c_1 - 1) + 1, 2t + 1)$
found in the lemma.

\medskip\noindent
Induction step. Say $0 \leq s < t$ we have $N, c$ as in $H_{s + 1}$.
Consider a module $M$ with $\ell(M) = s$. Then we can choose an $i$
such that $a_i^{c_1}$ is nonzero on $M$. It follows that
$\ell(M[a_i^c]) \geq s + 1$ and $\ell(M/a_i^{c_1}M) \geq s + 1$
and the induction hypothesis applies to them.
Consider the exact sequence
$$
0 \to M[a_i^{c_1}] \to M \xrightarrow{a_i^{c_1}} M \to M/a_i^{c_1}M \to 0
$$
Denote $E \subset M$ the image of the middle arrow.
Consider the corresponding diagram of Tor modules
$$
\xymatrix{
&
&
\text{Tor}_{i + 1}(M/a_i^{c_1}M, A/I^q) \ar[d] \\
\text{Tor}_i(M[a_i^{c_1}], A/I^q) \ar[r] &
\text{Tor}_i(M, A/I^q) \ar[r] \ar[rd]^0 &
\text{Tor}_i(E, A/I^q) \ar[d] \\
&
&
\text{Tor}_i(M, A/I^q)
}
$$
with exact rows and columns (for every $q$). The south-east arrow
is zero by our choice of $c_1$. We conclude that the
module $\text{Tor}_i(M, A/I^q)$ is sandwiched between
a quotient module of $\text{Tor}_i(M[a_i^{c_1}], A/I^q)$
and a submodule of $\text{Tor}_{i + 1}(M/a_i^{c_1}M, A/I^q)$.
Hence we conclude $H_s$ holds with $N$ replaced by $N + c$
and $c$ replaced by $2c$. Some details omitted.
\end{proof}

\end{document}
